在区分隔离状态的模型中，选择非隔离感染者作为控制对象有直接的传播学依据。Zhou 等\cite{Zhou2024SPCI}指出，隔离区外的传播主要由尚未被隔离或追踪的感染者造成，其中包括不易被发现的无症状感染者。在本文模型中，进入非隔离感染仓室 $I$ 和隔离感染仓室 $I_q$ 的全部新增感染均由社区中的 $I$ 引起。因此，在接触率有界的条件下，限制 $I$ 不仅控制社区感染峰值，也为两类病例的总新增感染流量提供统一上界。隔离感染者虽不再参与模型中的传播，仍可能需要救治，不能将其医疗需求排除在外。Zhang 等\cite{Zhang2026Behavior}在含住院过程的行为--传播模型中，也使用感染现患率上界讨论医疗压力显现之前的干预。

本文以避免医疗需求过载为现实动机，以社区非隔离感染峰值为直接控制对象。第\ref{sec:model:strategy}节通过全部新增感染流量，说明在共同的医疗需求关系和初始负担条件下，怎样选取足够低的社区感染阈值，使平均 ICU 需求不超过可用容量。这一联系不增加传播仓室，也不将恒定社区感染人数等同于 ICU 满负荷运行。后续分析以给定、可实施的阈值为基础，研究规定三阶段策略的解析结构和两种措施的分配；具体地区的临床容量校准与这一控制分析分别处理。
